%!TEX root = ../student_thesis.tex

\chapter{Introduction} \label{ch:introduction}
\section{Motivation} \label{sec:motivation}
% machine + converter
% the multirate problem –> monolithic, cost of FEM and why not
Most parts of a drive system, such as sources, switches, passive elements, and control circuitry, do not need to be modelled with the full spatial detail of Maxwell's partial differential equations \cite{Simonyi1979} and can be adequately described by lumped elements. These reduced models are formulated as differential-algebraic equations (DAEs), e.g. by modified nodal analysis (MNA) \cite{Chung-Wen1975}. For the machine itself, a lumped representation may be insufficient whenever saturation, eddy-current effects or the spatial distribution of the field govern its behaviour. A spatially distributed model is used instead, commonly obtained from the magnetoquasistatic (MQS) approximation of Maxwell's equations and discretised in space, e.g. by the finite element method (FEM) \cite{Monk2003}. Coupling such a field model to the circuit yields a system of partial differential-algebraic equations (PDAEs), which reduces to a system of DAEs once the field model is discretised in space \cite{Garcia2020}. Solving the system \textit{monolithically} would require a single time integrator to resolve the dynamics of the most active domain, so the latent domain is advanced with a step size far smaller than its dynamics require. Partitioning removes this restriction, as each subsystem can be integrated with a step size adapted to its own time constants. This is particularly beneficial when the time scales of the domains are separated by orders of magnitude, where it has been shown to reduce the computational effort considerably \cite{Schoeps2010}. Electric machines are increasingly fed by fast-switching power electronics, enabled by wide-bandgap semiconductor devices and driven by \textquote{ambitious} power-density and voltage targets for the coming decade, see Tab.~\ref{tab:doe-targets}. In such systems, the machine is the latent subsystem, as its electromagnetic and mechanical time constants exceed the converter switching period by orders of magnitude, which motivates multirate and co-simulation approaches \cite{SchoepsDiss, Pels2019}. Faster switching also increases the voltage slew rate $\mathrm{d}v/\mathrm{d}t$ at the machine terminals. Whereas insulation systems of standard machines are designed for slew rates of a few \si{\kilo\volt\per\micro\second}, silicon-carbide devices reach values exceeding \qty{100}{\kilo\volt\per\micro\second}, according to \cite{Petri2022}. Consequently, the behaviour of the machine at its terminals is no longer adequately represented by a lumped element, motivating a partitioned field-circuit approach in which the interface conditions capture the machine's terminal response.

Practically speaking, there are well-established solvers for each domain. However, commercial software is often closed-source and access to the internal variables during simulation is limited, making coupling the field-circuit domains difficult \cite{Gomes2017}. Hence, a so-called \textit{black-box} view of the problem, which is agnostic to the solvers, is the versatile implementation. More precisely, one might call it a \textit{grey-box} view, since we require some knowledge of the state of the field domain in the interface condition, as we will discuss later.

\begin{table}[htb]
  \centering
  \caption{Power electronics and electric motor targets (US DoE).}
  \label{tab:doe-targets}
  \begin{tabular}{lccc}
    \toprule
    & 2025 & 2030 & 2035 \\
    \midrule
    Voltage (V)                        	& 600 & 800 & 800   \\
    Power Density Electronics (kW/L)		& 150 & 200 & 225   \\
    Power Density Machines (kW/L)     	&  75 & 100 & 112.5 \\
    \bottomrule
  \end{tabular}
  \par\smallskip
  \footnotesize Data from \cite{USDRIVE_EDTT}.
\end{table}

\section{State of the Art}
% (monolithic solves, tight coupling / two-level Newton à la Xyce GenExt), classic waveform relaxation, ROM-based transmission conditions
% „Wo steht die Forschung heute, und warum braucht es meine Arbeit?“
Solving the system of DAEs is a thoroughly studied problem in the context of simulating field-circuit domains. Approaches differ in how tightly the subsystems are coupled, ranging from monolithic solves, where all equations are assembled into a single system, to loosely coupled \textit{dynamic iteration} (also called \textit{Waveform Relaxation} (WR)) where the subsystems are solved separately and exchange waveforms over time windows, see Section~\ref{sec:dynamic-iteration}. The monolithic approach requires a joint solver for both domains, which conflicts with the black-box setting and is prohibitively expensive in computation for the targeted application here, see Section~\ref{sec:motivation}. In the case of dynamic iteration, it has been identified that the exchange of algebraic variables can be a mathematical cause for divergence \cite{SchoepsDiss}. Thus, enhanced modelling of the interface via a \textit{surrogate} or \textit{reduced order model} (ROM) has been shown to improve convergence significantly. Without it, the iteration may be divergent \cite{Bartel2013}. While such an interface condition is straightforward to formulate analytically, its time discretisation involves the time integrators of both subsystems. When black-box tools are coupled, each subsystem discretises the shared condition with its own integrator, so the two discrete operators may differ.
Different operators result in an inconsistent interface condition and may possibly lead to a non-converging iteration, as we will see in Chapter~\ref{ch:benchmark}. This motivates the need for a consistent implementation of the interface without access to the solver's time integrator, see Chapter~\ref{ch:interface-discretisation}.
Xyce\texttrademark \ also provides a native coupling interface (General External Interface). Why it is not used for the dynamic iteration here is discussed in Section~\ref{sec:xyce-genext}.

\section{Objective and Contributions} \label{sec:objective-contributions}
% what's missing in Literature and what is the research issue (Forschungsfrage); explicit list of contributions of this thesis
The contributions of this thesis are:
\begin{enumerate}
	\item a dynamic iteration co-simulation of the field-circuit built entirely from black-box, netlist-driven simulators (demonstrated with Xyce\texttrademark);
	\item a consistent formulation and implementation of the interface condition that does not require access to either solver's time integrator;
	\item a study of the resulting scheme's convergence behaviour, transmission conditions and parameter-dependence, e.g. window-length, isolated from any particular field discretisation.
\end{enumerate}
Xyce\texttrademark \ is an open-source analogue circuit simulator used in this thesis. It can be replaced by other netlist-driven simulators with little change to the implementation, see Section~\ref{sec:xyce-why}.

\section{Outline of the Thesis}
% chapter one talks about ... etc.
In Chapter~\ref{ch:fundamentals}, the fundamental theory that is required to explain the implementation and discuss any results is laid out. In Chapter~\ref{ch:interface-discretisation} we analyse different discretisations of the interface condition. Chapter~\ref{ch:xyce} documents the implementation of the co-simulation with the Xyce\texttrademark \ parallel electronic simulator. In Chapter~\ref{ch:benchmark}, the implementation is verified against a suite of benchmark problems. Chapter~\ref{ch:application} contains parameter studies and a general approach to choosing them. Chapter~\ref{ch:conclusions} concludes the thesis and outlines possibilities for future work.